\documentclass[10pt,a4paper]{amsart}
\usepackage[margin=19mm]{geometry}
\usepackage{amsmath,amssymb,mathtools}
\usepackage[scr=rsfs,cal=euler]{mathalfa}
\usepackage{xfrac}
\usepackage[unicode,hidelinks,bookmarks=false]{hyperref}
\newcommand{\ie}{i.e.\ }
\newcommand{\cf}{cf.\ }
\newcommand{\resp}{respectively\ }
\newcommand{\Subsection}{\subsection}
\providecommand{\nobreakdash}{\nobreak\hbox{-}}
\setlength{\emergencystretch}{2em}
\multlinegap0pt
\usepackage{bm}
\usepackage{bbm}
\usepackage{dsfont}
\usepackage{xspace}
\usepackage{cancel}
\usepackage{mathtools}
\usepackage{amsthm}
\makeatletter
\@ifundefined{theorem}{\newtheorem{theorem}{Theorem}}{}
\@ifundefined{lemma}{\newtheorem{lemma}[theorem]{Lemma}}{}
\makeatother
\newcommand{\prob}{\mathbb{P}} %Probability
\title[On the size of temporal cliques in subcritical random tempor]{On the size of temporal cliques in subcritical random temporal graphs}
\author{Caelan Atamanchuk, Luc Devroye, Gabor Lugosi}
\date{Source: arXiv:2404.04462v3 (CC BY 4.0)}
\begin{document}
\maketitle

\noindent\emph{Source and licence.} On the size of temporal cliques in subcritical random temporal graphs. Version arXiv:2404.04462v3 (CC BY 4.0), distributed under the Creative Commons Attribution 4.0 International licence, as recorded in the arXiv licence metadata for this exact version and shown on the abstract page https://arxiv.org/abs/2404.04462v3. Full rights notice: licenses/arxiv-2404.04462-CC-BY-4.0.txt.\par\smallskip

\noindent\textit{Editorial scope.} Counted unit: the Lemma whose source label is \texttt{treesurvival} in the article (source0.tex lines 217--223, source label \texttt{treesurvival}), delivered together with the article's own complete original proof of it (source0.tex lines 226--268, one \texttt{proof} environment of 43 lines). No mathematics is written, repaired or re-derived: statement and proof are the article's own text line for line. Required context retained uncounted: treesep (lines 206--209). Every definition, display and label of those lines is retained unchanged. Omitted from this package and not redistributed: 1--205, 210--216, 224--225, 269--470. Nothing inside a retained excerpt is omitted or altered: every retained body is delivered exactly as the article's lines are written, so no omission receipt of any kind accompanies this package. Citations: the retained excerpts carry no citation command at all, so no bibliography entry of the article is transcribed and no reference is redistributed. Reference adaptation: none was needed and none was made; every reference inside the delivered unit resolves inside it, so no unresolved reference or citation is produced. Printed slips kept literal: no printed word, symbol, hypothesis or number of the counted statement or of its proof was altered. Layout and class: the article's own class is \texttt{article} with packages including geometry, setspace, fontenc, charter, euler, float, soul, bbm, verbatim, amsmath, amsthm, amssymb, textcomp, graphicx; the wrapper is the lane's pinned \texttt{amsart} editorial wrapper at 10pt on A4, which restates the theorem-like environments the retained text needs and adds the article's own macro definitions that the retained text needs (\texttt{\textbackslash prob}), with extra wrapper packages (bm, bbm, dsfont, xspace, cancel, mathtools). The delivered numbering is the unit's own. Assets: no retained span contains a figure environment, a graphics-inclusion command, a PDF-page inclusion, a table or a TikZ picture; no image file is copied into this package, the package declares no asset dependency and no image file is redistributed with it. Licence: version arXiv:2404.04462v3 of this article is distributed under the Creative Commons Attribution 4.0 International licence, as recorded in the arXiv licence metadata for this exact version and shown on the abstract page https://arxiv.org/abs/2404.04462v3. A full rights notice is delivered at licenses/arxiv-2404.04462-CC-BY-4.0.txt. Characters: every retained excerpt is pure ASCII, so the delivered engine, XeLaTeX with its default Latin Modern text font, reports no missing character.
\begin{lemma}\label{treesep}
Let $P_1,\ldots,P_q$ be a finite collection of distinct infinite paths in $T$, and let $(X_k)_{k \geq 0}$ be a random walk down the tree, that is, $X_0$ is the root and $X_k$ is uniformly distributed over the children of $X_{k-1}$ for all $k \geq 1$. Then, $\prob(\tau \geq \ell) \leq \frac{q}{n^\ell}$, where
$$\tau = \max\{ k > 0 : X_0,\ldots,X_k \text{ coincides with one of the $q$ paths}\}.$$
\end{lemma}

\begin{lemma}\label{treesurvival}
Let $T^*$ be the set of all vertices that are reachable from the root of $T$. If $v_1,\ldots,v_q$ are uniform vertices chosen from the $\ell$-th generation of $T$, then
$$\prob(v_1,\ldots,v_q \in T^*) \leq \frac{(q-1)!e^{npq}}{n^{\ell q}}.$$
Furthermore, when $\ell \geq (np)^4$ and $np \to \infty$ as $n \to \infty$,
$$\prob(v_1,\ldots,v_q \in T^*) = O\left(\frac{(q-1)!C^{q-1}}{n^{\ell q}}\right),$$
for some $C > 0$.
\end{lemma}

\begin{proof}
Suppose that $v_1,\ldots,v_r \in T^*$ and let $T'$ be a subtree of $T$ consisting of $r$ distinct infinite paths starting at the root through $v_1,\ldots,v_r$. Let $(X_k)_{k \geq 0}$ be a random walk down the tree (independent of $v_1,\ldots,v_r$) and let $\tau$ be as in Lemma \ref{treesep}. If $\tau = j$, then there are $\ell-j$ edges left that need to both exist and be increasing to have $X_\ell \in T^*$. Hence, if $V_r = \{v_1 , \ldots , v_r \in T^*\}$,
$$\prob(X_\ell \in T^* | V_r \cap \{\tau = j\}) = p^{\ell-j}\prob\big(\big\{\text{a path of length $\ell$ is increasing}\big\} \big|\big\{ \text{first $j$ edges are increasing}\big\}\big) = \frac{p^{\ell-j}j!}{\ell!}.$$
Combining this with Lemma \ref{treesep} yields
\begin{align}\label{sum}
\prob(X_\ell \in T^* | V_r)  &
\leq \sum_{j=0}^{\ell-1} \prob(\tau = j)\prob(X_\ell \in T^* | V_r \cap \{\tau = j\}) 
+ \prob(\tau \ge \ell)\prob(X_\ell \in T^* | V_r \cap \{\tau \ge \ell\}) 
 \nonumber \\
& \leq \sum_{j=0}^\ell \left(\frac{r}{n^j}\right)\frac{p^{\ell-j}j!}{\ell!} = \frac{r}{n^\ell}\sum_{j=0}^\ell\frac{(np)^{\ell-j}j!}{\ell!}.
\end{align}
To get the first bound we use the fact that $\frac{j!}{\ell!} \leq \frac{1}{(\ell-j)!}$ to get 
$$\prob(X_\ell \in T^*|V_r) \leq \frac{re^{np}}{n^\ell}.$$
Then, since $X_\ell$ is distributed uniformly across the $\ell$-th generation,  applying the above inequality repeatedly,
\begin{align*}
\prob(v_1,\ldots,v_q \in T^*) &= \prob(v_1 \in T^*) \cdot \prob(v_2 \in T^* | V_1) \cdots \prob(v_q \in T^* | V_{q-1}) \\
&\leq \frac{p^\ell}{\ell!}\prod_{r=1}^{q-1}\left(\frac{re^{np}}{n^\ell}\right) = \left(\frac{(np)^{\ell}}{\ell!}\right)\frac{(q-1)!e^{np(q-1)}}{n^{\ell q}} \leq \frac{(q-1)!e^{npq}}{n^{\ell q}}.
\end{align*} 
For the second inequality, we split the sum in (\ref{sum}) in two separate pieces
\begin{align*}
    A = \sum_{j=0}^{\ell - \sqrt{\ell}}\frac{(np)^{\ell - j}j!}{\ell!}, \quad \text{and} \quad B = \sum_{j=\ell -\sqrt{\ell}+1}^\ell \frac{(np)^{\ell-j}j!}{\ell!} = \sum_{k=0}^{\sqrt{\ell}-1} \frac{(np)^k}{\ell \cdot (\ell-1) \cdots (\ell-k+1)}.
\end{align*}
The first term may be bounded by
\begin{align*}
A &\leq \sum_{k=\sqrt{\ell}}^\ell \frac{(np)^k}{k!} \leq \left|e^{np}- \sum_{k=0}^{\sqrt{\ell}-1} \frac{(np)^k}{k!}\right| \leq \frac{e^{np}(np)^{\sqrt{\ell}}}{(\sqrt{\ell})!}~,
\end{align*}
where the second inequality follows from the Lagrange form of the remainder in Taylor's theorem.

For the second term, since $\ell \cdot (\ell-1) \cdots (\ell - k + 1) \geq \ell^{k}(1-\frac{1}{\sqrt{\ell}})^{\sqrt{\ell}}$ for all $0 \leq k \leq \sqrt{\ell}$, we have that
\begin{align*}
    B &\leq \frac{1}{\left(1-\frac{1}{\sqrt{\ell}}\right)^{\sqrt{\ell}}}\sum_{k=0}^{\sqrt{\ell}-1} \left(\frac{np}{\ell}\right)^k \leq \frac{1}{\left(1-\frac{np}{\ell}\right)\left(1-\frac{1}{\sqrt{\ell}}\right)^{\sqrt{\ell}}} = O(1),
\end{align*}
when $\ell \geq (np)^4$. Combining both the bounds along with Stirling's approximation, we conclude that there is some $C > 0$ such that
$$
\prob(X_\ell \in T^* | V_r) 
\leq \frac{r}{n^\ell}\left(\frac{e^{np}(np)^{\sqrt{\ell}}}{\sqrt{\ell}!} + O(1)\right) \leq \frac{r}{n^\ell}\left(\frac{e^{np}(np)^{(np)^2}e^{(np)^2}}{(np)^{2(np)^2}}+O(1)\right) \leq \frac{Cr}{n^\ell},$$
when $\ell \geq (np)^4$ and $np \to \infty$. Proceeding exactly as we did for the first inequality,
\begin{align*}
\prob(v_1,\ldots,v_q \in T^*) &= \prob(v_1 \in T^*) \cdot \prob(v_2 \in T^* | V_1) \cdots \prob(v_q \in T^* | V_{q-1}) \\
&\leq \frac{p^\ell}{\ell!}\frac{(q-1)!C^{q-1}}{n^{\ell (q-1)}} \leq \frac{(np)^{(np)^4}}{(np)^4!}\frac{(q-1)!C^{q-1}}{n^{\ell q}} = O\left( \frac{(q-1)!C^{q-1}}{n^{\ell q}}\right),
\end{align*} 
where in the final bound we use the fact that $\frac{x^{x^4}}{(x^4)!} = o(1)$ as $x\to \infty$, which is an immediate consequence of Stirling's approximation.
\end{proof}

\end{document}
